\section{Exact Local Demand Calculus}
\label{sec:exact-local-demand-calculus}

This section develops the local calculus for configurations in which a coarse
length sum does not by itself record enough information.  The mechanism is
always the same.  A center trace fixes boundary inputs and complementary
radial demands; an exact
local map propagates these data through nonsupercritical V triangles; and the
returning demand exceeds the remaining boundary or radial capacity.

Actual maximal reaches continue to be denoted by $(A_i,B_i,C_i)$, while
lowercase letters denote prescribed demands.  Every estimate in this section
is proved analytically.  In particular, no empirical branch catalogue,
machine-assisted certificate, or formal root from a discarded algebraic
component is used.

\subsection{The exact local feasible set}

Put
\[
 u=\left(\frac12,\frac{\sqrt3}{2}\right),\qquad
 v=\left(\frac12,-\frac{\sqrt3}{2}\right).
\]
For $0\le a,b,c\le1$ define
\[
 K(a,b,c)=\conv\{0,au,bv,c(u+v)\}.
\]
Thus $a,b$ are lower demands on the two incident boundary rays and $c$ is a
lower demand on the inward radial ray.

\begin{definition}
\label{def:admissible-demands}
The local admissible set is
\[
 \begin{aligned}
 \mathcal A=\{(a,b,c)\in[0,1]^3:\;&K(a,b,c)\text{ is contained}\\
 &\text{in a closed unit equilateral triangle}\}.
 \end{aligned}
\]
For $a,c\in[0,1]$ put
\[
 \begin{aligned}
 B_c(a)&=\max\{b:(a,b,c)\in\mathcal A\},\\
 F_c(a)&=\min\{B_c(a),1-a\},\\
 G_c(a)&=1-F_c(a).
 \end{aligned}
\]
\end{definition}

The distinction between demands and actual reaches is important.  If a role
has actual reaches $(A,B,C)$ and $A\ge a$, $C\ge c$, then $(a,B,c)$ is a
demand triple realized by the same role.  If in addition $A+B\le1$, then
\begin{equation}
 B\le F_c(a).
 \label{eq:safe-capped-bound}
\end{equation}
No assertion about the type of a maximizing triangle is required.

\begin{proposition}[Exact admissible set and radial envelope]
\label{prop:exact-admissible-set}
\label{prop:local-admissible}
Put
\[
 s=a+b,\qquad d=a^2+ab+b^2,\qquad
 m=\min\{a,b\},\qquad M=\max\{a,b\},
\]
and
\[
 q=s^4-s^2+ab.
\]
Define
\begin{align*}
 F_L(a,b,c)&=c^4-c^2+mc-m^2,\\
 F_T(a,b,c)&=(s^2-1)c^2+Mc-M^2,\\
 F_S(a,b,c)&=(m^2-1)c^2+(2mM^2+M)c+M^4-M^2.
\end{align*}
Then $\mathcal A=\mathcal A_L\cup\mathcal A_T\cup\mathcal A_S$, where
\begin{align*}
 \mathcal A_L&=\{d\le1, s\le1, q\le0, F_L\le0\},\\
 \mathcal A_T&=\{d\le1, s\le1, q\ge0, F_T\le0, c\le2M\},\\
 \mathcal A_S&=\{d\le1, s>1, F_S\le0, c\le\tfrac12\}.
\end{align*}
All variables in these three sets are also constrained to $[0,1]$.

For $d\le1$, the maximal radial demand is
\begin{equation}
 c_{\max}(a,b)=
 \begin{cases}
  1,&a=b=0,\\
  C_L(m),&s\le1, q\le0, (a,b)\ne(0,0),\\
  C_T(m,M),&s\le1, q>0,\\
  C_S(m,M),&s>1,
 \end{cases}
 \label{eq:radial-envelope}
\end{equation}
where $C_L(m)$ is the unique root in $[\sqrt3/2,1]$ of
\[
 z^4-z^2+mz-m^2=0,
\]
and
\[
 C_T(m,M)=\frac{2M}{1+\sqrt{4s^2-3}},
\qquad
 C_S(m,M)=
 \frac{M(2mM+1)-M\sqrt{4d-3}}{2(1-m^2)}.
\]
If $d>1$, the radial envelope is undefined.  At $q=0$ the two subcritical
formulas agree and equal $s$.
\end{proposition}

\begin{proof}
If $s>0$ and $cs\le ab$, then $c/a+c/b\le1$ (with the zero-coordinate
cases obtained by continuity), so $c(u+v)$ lies in
$\conv\{0,au,bv\}$.  Its minimum enclosing equilateral side is the diameter
\[
 \sqrt{a^2+ab+b^2}.
\]
Assume now that $cs>ab$.  The four points occur cyclically as
$0,bv,c(u+v),au$.  For three outward unit normals separated by
$120^\circ$, the required side equals $2/\sqrt3$ times the sum of the three
support values.  On an angular cell with fixed support points this sum has
second derivative equal to its negative.  Hence an angular minimum occurs
when a normal is perpendicular to one of the four hull edges.  The four
resulting side lengths are
\begin{align*}
 L_{0A}&=b+\max\{a,c\},&
 L_{B0}&=a+\max\{b,c\},\\
 L_{AC}&=
 \begin{cases}
 (a^2+bc)/\sqrt{a^2-ac+c^2},&a\ge c,\\
 c(a+b)/\sqrt{a^2-ac+c^2},&a<c\le a+b,\\
 c^2/\sqrt{a^2-ac+c^2},&c>a+b,
 \end{cases}
\end{align*}
and $L_{CB}$ is obtained by interchanging $a,b$.  Thus the minimum side is
the minimum of these four expressions.

Ordering $a,b$ as $m\le M$, squaring the positive support equations gives
$F_L=0,F_T=0,F_S=0$ in the three respective contact regimes; the inactive
support inequalities give $q\le0,q\ge0,s>1$.  Squaring creates one spurious
component in each of the last two cells.  In the $T$ cell the two roots are
\[
 \frac{2M}{1+\sqrt{4s^2-3}},
 \qquad
 \frac{2M}{1-\sqrt{4s^2-3}},
\]
and $c\le2M$ selects the first.  In the $S$ cell the smaller root lies below
$1/2$ and the other above $1/2$; hence $c\le1/2$ is exactly the geometric
selector.  The $L$ fiber is the interval from zero to its unique selected
root.  Solving the selected equations gives \eqref{eq:radial-envelope}.
This also proves that $\mathcal A$ is coordinatewise down-closed.
\end{proof}

\subsection{The exact capped map}

For $1/2<c<1$ put
\[
 \ell(c)=\frac c2\left(1-\sqrt{4c^2-3}\right),\qquad
 r(c)=c-\ell(c)
\]
when $c\ge\sqrt3/2$, and
use the order-transition function $h(c)$ from
\eqref{eq:signed-order-transition}.
Also define
\begin{align*}
 T_-(a,c)&=\frac{\sqrt{a^2-ac+c^2}}c-a,\\
 T_+(a,c)&=
 \frac{2(1-a^2)c^2}
 {2ac^2+c+\sqrt{(2ac^2+c)^2-4(1-c^2)(1-a^2)c^2}}.
\end{align*}

\begin{proposition}[Four-label capped map]
\label{prop:capped-demand-map}
For $1/2<c\le\sqrt3/2$,
\begin{equation}
 F_c(a)=
 \begin{cases}
  1-a,&0\le a\le1-c,\\
  T_+(a,c),&1-c<a<h(c),\\
  h(c),&a=h(c),\\
  T_-(a,c),&h(c)<a<c,\\
  1-a,&c\le a\le1.
 \end{cases}
 \label{eq:S2-capped-map-low}
\end{equation}
For $\sqrt3/2<c<1$,
\begin{equation}
 F_c(a)=
 \begin{cases}
  1-a,&0\le a\le1-c,\\
  T_+(a,c),&1-c<a\le\ell(c),\\
  \ell(c),&\ell(c)<a<r(c),\\
  T_-(a,c),&r(c)\le a<c,\\
  1-a,&c\le a\le1.
 \end{cases}
 \label{eq:S2-capped-map-high}
\end{equation}
Thus there are exactly four labels
\[
 \mathrm{Full},\qquad L,\qquad T_-,\qquad T_+^{\mathrm{hi}}.
\]
At $a=\ell(c)$ in \eqref{eq:S2-capped-map-high},
$F_c(a)=r(c)$; immediately to its right the value is $\ell(c)$.  This
downward jump is genuine.  The formal other root of the $T_+$ quadratic is
not feasible because it violates $c\le2\max\{a,F_c(a)\}$.

The maps $F_c$ are nonincreasing, the maps $G_c=1-F_c$ are nondecreasing
and extensive, and
\begin{equation}
 G_c(a)\le z\quad\Longleftrightarrow\quad
 G_c(1-z)\le1-a.
 \label{eq:S2-capped-duality}
\end{equation}
\end{proposition}

\begin{proof}
The cap $a+b\le1$ eliminates the $S$ cell of
Proposition~\ref{prop:exact-admissible-set}.  On the cap $b=1-a$, the $T$
inequality reduces to $M(c-M)\le0$; hence the Full face is feasible exactly
for $a\le1-c$ or $a\ge c$.  Off that face a maximal point lies on $F_L=0$
or $F_T=0$.  The $L$ equation has roots $\ell(c),r(c)$ and the selector
$q\le0$ keeps only the smaller coordinate.  The ordered $T$ halves give
respectively the displayed lower quadratic root $T_+$ and the root $T_-$.
Their order boundary is $a=b=h(c)$, while their $q=0$ boundaries are
$(\ell(c),r(c))$ and $(r(c),\ell(c))$.  These observations give
\eqref{eq:S2-capped-map-low}--\eqref{eq:S2-capped-map-high}, including all equality
assignments.

Down-closedness gives monotonicity, while $F_c(a)\le1-a$ gives
$G_c(a)\ge a$.  Finally,
\[
 G_c(a)\le z\iff F_c(a)\ge1-z
 \iff(a,1-z,c)\in\mathcal A,\quad a+1-z\le1.
\]
Reflection $a\leftrightarrow b$ gives \eqref{eq:S2-capped-duality}.
\end{proof}

Put
\[
 \mu=\frac{2\sqrt3-3}{4},
\]
and, for $0<X<1-\sqrt3/2$, write
\[
 e(X)=\ell(1-X).
\]
The root equation gives
\begin{equation}
 \begin{aligned}
 X&<e(X) &&\left(0<X<1-\frac{\sqrt3}{2}\right),\\
 e(X)&<\frac{2X}{1-2X} &&(0<X\le\mu),\\
 e(X)&\le2X+5X^2 &&(0<X\le1/8).
 \end{aligned}
 \label{eq:low-root-bounds}
\end{equation}
Moreover,
\begin{equation}
 a>e(X)\quad\Longrightarrow\quad
 G_{1-X}(a)\ge1-e(X)
 \qquad\left(0<X<1-\frac{\sqrt3}{2}\right).
 \label{eq:high-demand-threshold}
\end{equation}
Indeed, for
\[
 p_X(z)=z^2-(1-X)z+(1-X)^2-(1-X)^4,
\]
one has
\[
 p_X(X)=X(1-3X+4X^2-X^3)>0.
\]
Since $X<(1-X)/2$, this puts $X$ strictly below the smaller root $e(X)$.
For $U=2X/(1-2X)$, direct substitution gives
\[
 p_X(U)=-\frac{X^2}{(1-2X)^2}
 (4X^4-20X^3+37X^2-28X+3).
\]
The polynomial in parentheses is decreasing on $[0,\mu]$ and at the right
endpoint equals $8217/64-74\sqrt3>0$.  Indeed, its derivative is at most
$16/512+74/8-28<0$, and
$(8217/64)^2-3\mathbin{\cdot}74^2=230001/4096>0$.  Hence $p_X(U)<0$, which puts $U$
strictly between the roots.  Finally,
\[
 p_X(2X+5X^2)=X^2(3X+4)(8X-1)\le0,
\]
and $X<2X+5X^2<(1-X)/2$ on the stated range.  Thus the trial point lies
between the two roots, proving the third comparison.  The implication
\eqref{eq:high-demand-threshold} is read directly from
\eqref{eq:S2-capped-map-high}.  We shall also use
$e(X)<2X+3X^2$ for $0<X\le1/10$, obtained from
$p_X(2X+3X^2)=X^2(8X^2+19X-2)<0$; here too the trial point lies strictly
between $X$ and $(1-X)/2$.
